\documentclass[11pt]{amsart}

\usepackage[T1]{fontenc}
\usepackage{lmodern}
\usepackage{microtype}
\usepackage{amsmath,amssymb}
\usepackage[hidelinks]{hyperref}

\hypersetup{
  pdftitle={Counterexamples to the Glushkov MFN Conjectures}
}

\newtheorem{theorem}{Theorem}
\newtheorem{corollary}[theorem]{Corollary}

\newcommand{\Gl}{\operatorname{Gl}}
\newcommand{\MFN}{\operatorname{MFN}}
\newcommand{\IARp}{\operatorname{IAR}^{+}}

\title[Glushkov MFN counterexamples]
{Counterexamples to the Glushkov MFN Conjectures}

\date{}

\subjclass[2020]{68Q45}
\keywords{regular expressions, Glushkov automata, memoization, ambiguity,
feedback vertex set}

\begin{document}

\begin{abstract}
We refute the Glushkov parts of Conjectures~1 and~2 of Berglund,
van der Merwe, and le Roux. For every $k\ge 2$, we give a regular
expression for which the Minimum Feedback Node (MFN) set of the
Glushkov automaton has size $k$, whereas the unique minimum
memoization set is a singleton outside MFN. We characterize all
memoization sets eliminating infinite degree of ambiguity and show
that every ordering of $\IARp$ returns the entire MFN set. Both
output-to-optimum ratios therefore equal $k$ and are unbounded. The
corresponding Thompson questions are not addressed.
\end{abstract}

\maketitle

\section{The counterexample}

Berglund, van der Merwe, and le Roux conjectured that, for every
regular expression, a minimum memoization set can be chosen inside the
Minimum Feedback Node (MFN) set, for both the Thompson and Glushkov
constructions \cite[Conjecture~1]{BerglundEtAl2026}. They equivalently
conjectured that some ordering makes $\IARp$ return a minimum
memoization set
\cite[Conjecture~2 and Observation~3]{BerglundEtAl2026}.
We disprove both Glushkov assertions.

We use the conventions of \cite{BerglundEtAl2026}, in particular that
$+$ is a primitive regular-expression operator. For an NFA $A$ and
$M\subseteq Q(A)$, write $A[M]$ for the corresponding memoized NFA.
By the IDA characterization in \cite[Theorem~1]{BerglundEtAl2026},
$A[M]$ has infinite degree of ambiguity (IDA) if and only if there are
distinct states $p,q$ and a nonempty word $v$ such that
\[
p\xrightarrow{v}p,
\qquad
p\xrightarrow{v}q,
\qquad
q\xrightarrow{v}q,
\]
with the $q$-cycle containing no state of $M$. Define
\[
\mu(A)=
\min\bigl\{
  |M|:M\subseteq Q(A),\ A[M]\text{ has no IDA}
\bigr\}.
\]
For an ordering $\prec$ of $\MFN(A)$, let $\IARp_{\prec}(A)$ denote
the output of the procedure in
\cite[Definition~5]{BerglundEtAl2026}.

For $k\ge 2$, let
\[
R_k=
\Bigl(
  (\underbrace{a^+\mid\cdots\mid a^+}_{k\text{ occurrences}})b
\Bigr)^*,
\qquad
A_k=\Gl(R_k).
\]
The alternatives are distinct syntax-tree occurrences and hence
distinct positions in the Glushkov linearization; no algebraic
simplification is applied before the construction. Write $a_i^+$ for
the $i$th alternative. Let $x_1,\ldots,x_k$ be the positions labelled
$a$, let $y$ be the position labelled $b$, and put
\[
X_k=\{x_1,\ldots,x_k\}.
\]
Besides the initial state $q_0$, the nonempty transitions are
\[
\delta(q_0,a)=X_k,
\qquad
\delta(x_i,a)=\{x_i\},
\qquad
\delta(x_i,b)=\{y\},
\qquad
\delta(y,a)=X_k
\]
for $1\le i\le k$. The final states are $q_0$ and $y$, and all states
are useful.

\begin{theorem}
For every $k\ge 2$ and every $M\subseteq Q(A_k)$,
\[
A_k[M]\text{ has no IDA}
\quad\Longleftrightarrow\quad
y\in M\ \text{ or }\ X_k\subseteq M.
\]
Moreover,
\[
\MFN(A_k)=X_k,
\qquad
\mu(A_k)=1,
\]
and $\{y\}$ is the unique minimum memoization set. For every ordering
$\prec$ of $X_k$,
\[
\IARp_{\prec}(A_k)=X_k.
\]
\end{theorem}

\begin{proof}
Algorithm~2 of \cite{BerglundEtAl2026} gives
\[
a_i^+
\longmapsto
(\varnothing,\{x_i\},\mathrm{false}),
\qquad
(a_1^+\mid\cdots\mid a_k^+)
\longmapsto
(\varnothing,X_k,\mathrm{false}).
\]
Concatenation with $b$ preserves the latter tuple, and the outer star
changes only its nullability. Thus the root tuple is
\[
(R_{\mathrm{op}},R_{\mathrm{req}},
  \operatorname{nullable}(R_k))
=
(\varnothing,X_k,\mathrm{true}),
\]
so
\[
\MFN(A_k)=X_k.
\]
Independently, every feedback vertex set contains all of $X_k$, since
each $x_i$ has a self-loop, while deleting $X_k$ leaves no directed
cycle.

Suppose first that $y\in M$. Every directed cycle avoiding $M$ is then
a self-loop at some $x_i$, and every word labelling such a cycle is
$a^r$ for some $r\ge 1$. If a state $p$ returns to itself on $a^r$,
then $p=x_j$ for some $j$, and
\[
\delta(x_j,a^r)=\{x_j\}.
\]
Thus
\[
p\xrightarrow{a^r}q
\qquad\text{and}\qquad
q\xrightarrow{a^r}q
\]
force $p=q$, so the IDA criterion cannot hold. If instead
$X_k\subseteq M$, every directed cycle meets $M$, and again there is
no IDA.

Conversely, suppose that
\[
y\notin M
\qquad\text{and}\qquad
X_k\nsubseteq M.
\]
Choose $x_j\in X_k\setminus M$ and an index $\ell\ne j$. With
\[
p=x_\ell,
\qquad
q=x_j,
\qquad
v=ba,
\]
we have
\[
x_\ell\xrightarrow{ba}x_\ell,
\qquad
x_\ell\xrightarrow{ba}x_j,
\qquad
x_j\xrightarrow{ba}x_j.
\]
The $q$-cycle visits only $x_j$ and $y$, neither of which lies in $M$.
The IDA criterion imposes no avoidance condition on the other two
paths, so $x_\ell$ may itself lie in $M$. Hence $A_k[M]$ has IDA.
This proves the characterization.

The empty set therefore fails, while $\{y\}$ succeeds, so
\[
\mu(A_k)=1.
\]
The characterization also shows that every successful singleton is
$\{y\}$, proving uniqueness.

Finally, $\IARp$ starts from $X_k$. Removing any $x_i$ produces
$X_k\setminus\{x_i\}$, which contains neither $y$ nor all of $X_k$ and
therefore has IDA. Thus no state is removable, regardless of the
ordering.
\end{proof}

\begin{corollary}
The Glushkov parts of Conjectures~1 and~2 in
\cite{BerglundEtAl2026} are false. In fact, for every $k\ge 2$ and
every ordering $\prec$,
\[
\frac{|\MFN(A_k)|}{\mu(A_k)}
=
\frac{|\IARp_{\prec}(A_k)|}{\mu(A_k)}
=
k.
\]
Thus both gaps are unbounded.
\end{corollary}

\begin{proof}
The unique minimum memoization set is $\{y\}$, whereas
\[
y\notin\MFN(A_k)=X_k.
\]
Consequently, no minimum memoization set is contained in MFN. The
theorem also shows that every ordering of $\IARp$ returns $X_k$. The
ratios follow from $|X_k|=k$ and $\mu(A_k)=1$.
\end{proof}

\subsection{A second family}

The same two failures occur with distinct terminal letters and a
minimum memoization set of size $2$. For $k\ge 2$, let
\[
S_k=
\Bigl(
  \bigl(
    a\mid\underbrace{a^+\mid\cdots\mid a^+}_{k\text{ occurrences}}
  \bigr)(b\mid c)
\Bigr)^*,
\qquad
G_k=\Gl(S_k),
\]
again constructed from the displayed syntax without algebraic
simplification, so that the repeated alternatives are distinct
positions. Let $z_0$ be the position of the plain $a$, let
$z_1,\ldots,z_k$ be the positions of the $k$ occurrences of $a^+$, let
$y_b$ and $y_c$ be the positions labelled $b$ and $c$, and put
$Z_k=\{z_0,z_1,\ldots,z_k\}$. Besides the initial state $q_0$, the
nonempty transitions are
\[
\delta(q_0,a)=\delta(y_b,a)=\delta(y_c,a)=Z_k,
\qquad
\delta(z_i,b)=\{y_b\},
\qquad
\delta(z_i,c)=\{y_c\}
\]
for $0\le i\le k$, together with $\delta(z_i,a)=\{z_i\}$ for
$1\le i\le k$. The final states are $q_0$, $y_b$ and $y_c$.

Here $Z_k$ is the unique minimum directed feedback vertex set of $G_k$:
each $z_i$ with $i\ge 1$ carries a self-loop and hence lies in every
feedback vertex set; deleting those $k$ states leaves the two cycles
$z_0\to y_b\to z_0$ and $z_0\to y_c\to z_0$, of which $z_0$ is the only
common state; and deleting $Z_k$ leaves $q_0,y_b,y_c$ with no
transitions among them. Since the MFN algorithm returns a minimum
feedback vertex set \cite[Theorem~4]{BerglundEtAl2026},
\[
\MFN(G_k)=Z_k,
\qquad
|\MFN(G_k)|=k+1.
\]
Memoizing $\{y_b,y_c\}$ eliminates IDA: every directed cycle avoiding
that set is a self-loop at some $z_j$ with $j\ge 1$, so its label is
$a^r$ with $r\ge 1$, and any state returning to itself on $a^r$ is such
a $z_j$; as $\delta(z_j,a^r)=\{z_j\}$ is a singleton, the two states of
the IDA criterion coincide. Conversely, suppose $M$ eliminates IDA and
$|M|\le 2$. If $y_b\notin M$, then, since $|Z_k|=k+1\ge 3$, we may
choose $q\in Z_k\setminus M$ and $p\in Z_k\setminus\{q\}$, and
\[
p\xrightarrow{ba}p,
\qquad
p\xrightarrow{ba}q,
\qquad
q\xrightarrow{ba}q,
\]
where the $q$-cycle visits only $q$ and $y_b$, neither in $M$; so
$y_b\in M$, and the word $ca$ gives $y_c\in M$ in the same way. Hence
\[
\mu(G_k)=2,
\]
and $\{y_b,y_c\}$ is the unique minimum memoization set. It is disjoint
from $\MFN(G_k)$. The same $ba$ witness applies to every proper subset
of $Z_k$, since $y_b\notin Z_k$, so no state of the initial MFN set is
removable and $\IARp_{\prec}(G_k)=Z_k$ for every ordering $\prec$. Thus
$S_k$ refutes the same two Glushkov assertions, with
\[
\frac{|\MFN(G_k)|}{\mu(G_k)}
=
\frac{|\IARp_{\prec}(G_k)|}{\mu(G_k)}
=
\frac{k+1}{2},
\]
which is again unbounded. For $k=2$ this is the five-position
expression $\bigl((a\mid a^+\mid a^+)(b\mid c)\bigr)^*$.

\begin{thebibliography}{9}

\bibitem{BerglundEtAl2026}
M.~Berglund, B.~van der Merwe, and I.~le Roux,
\emph{Selective memoization for efficient backtracking regular
expression matching},
Electron. Proc. Theor. Comput. Sci. \textbf{446} (2026), 1--17.
\href{https://doi.org/10.4204/EPTCS.446.1}
{doi:10.4204/EPTCS.446.1}.

\end{thebibliography}

\end{document}